% BHU PYQ -- Manually transcribed & error-corrected
\documentclass[11pt,a4paper]{article}

\usepackage[a4paper,top=2.5cm,bottom=2.5cm,left=2.8cm,right=2.8cm,headheight=14pt]{geometry}
\usepackage{amsmath,amssymb}
\usepackage[T1]{fontenc}
\usepackage[utf8]{inputenc}
\usepackage{lmodern}
\usepackage{microtype}
\usepackage{fancyhdr}
\usepackage{enumitem}
\usepackage{hyperref}
\usepackage{lastpage}
\usepackage{parskip}

\hypersetup{colorlinks=true,urlcolor=black,linkcolor=black,
  pdftitle={BPT-401: Electronics and Modern Physics -- BHU PYQ},pdfauthor={Banaras Hindu University}}

\pagestyle{fancy}\fancyhf{}
\renewcommand{\headrulewidth}{0.4pt}\renewcommand{\footrulewidth}{0.4pt}
\fancyhead[L]{\small   \textit{Banaras Hindu University}}
\fancyhead[R]{\small   \textit{BPT-401: Electronics and Modern Physics}}
\fancyfoot[C]{\small Page  \thepage\ of \pageref{LastPage}}


\newlist{parts}{enumerate}{2}
\setlist[parts,1]{label=(\alph*),leftmargin=*,itemsep=5pt,topsep=3pt}
\setlist[parts,2]{label=(\roman*),leftmargin=*,itemsep=3pt,topsep=2pt}

\newcommand{\pts}[1]{\hfill   \textit{[#1]}}


\begin{document}


\begin{center}
  {\large\bfseries BANARAS HINDU UNIVERSITY}\\[2pt]
  {
\normalsize B.Sc. (Hons.) Semester IV Examination, 2013-14}\\[6pt]
  \rule{0.8\linewidth}{0.5pt}\\[6pt]
  {\large\bfseries Paper No. BPT-401: Electronics and Modern Physics}\\[4pt]
  {
\normalsize Time: 3 Hours\hfill Maximum Marks: 70}
\end{center}
\rule{\linewidth}{0.4pt}
\smallskip



\noindent   \textit{Instructions: Answer   \textbf{five} questions including Question~1 which is compulsory. Symbols have their usual meanings.}
\bigskip




\noindent   \textbf{Question 1.} Answer any   \textbf{five} of the following: \pts{5 $ \times$ 4 = 20}

\begin{parts}
  \item What are polar and non-polar dielectrics? Explain the terms electronic, ionic, and dipolar polarizabilities.
  \item Electrons are accelerated in a TV tube through a potential difference of $10\, \text{kV}$. Find the highest frequency of the electromagnetic wave emitted when these electrons strike the screen of the tube. Justify that these waves are X-rays.
  \item Explain the various parameters in the given equation for a p-n junction diode:
    \[ I = I_0 \left[ e^{\frac{eV}{\eta k T}} - 1 \right] \]
    Prove that the dynamic resistance of the diode is given by $r_d = \frac{\eta k T}{e I}$.
  \item How is current gain $\alpha$ different from $\beta$? Explain with expressions. Determine the current gain $\beta$ and emitter current $I_E$ for a transistor circuit in CE configuration, where $I_B = 50\,\mu \text{A}$ and $I_C = 3.65\, \text{mA}$.
  \item Show that the de Broglie wavelength of a particle of mass $m$, charge $q$, and accelerated through a potential difference of $V$ volts is:
    \[ \lambda = \frac{h}{\sqrt{2mqV}} \]
  \item Explain the important features of paramagnetic, diamagnetic, and ferromagnetic materials. Give the domain theory of ferromagnetism.
  \item Explain why a transistor needs to be biased. Draw the load line and explain its importance.
\end{parts}

\medskip\rule{\linewidth}{0.2pt}




\noindent   \textbf{Question 2.}

\begin{parts}
  \item Explain the formation of the depletion region in a P-N junction diode. Find the expression for the width of the depletion region and the height of the potential barrier without external voltage. \pts{10}
  \item A silicon p-n junction has a depletion width of $20\,\mu \text{m}$ at a reverse voltage of $16\, \text{V}$. If the reverse voltage is increased to $25\, \text{V}$, find the new depletion width. \pts{4}
\end{parts}

\medskip\rule{\linewidth}{0.2pt}




\noindent   \textbf{Question 3.}

\begin{parts}
  \item What is Lorentz local field? Derive the Lorentz formula for the internal electric field. Hence, derive the Clausius-Mossotti equation:
    \[ \frac{K - 1}{K + 2} = \frac{N \alpha_e}{3 \varepsilon_0} \]
    where the symbols have their usual meanings. \pts{10}
  \item A slab of dielectric material of dielectric constant $10$ and thickness $5\, \text{mm}$ is placed between the parallel plates of a capacitor. If a potential difference of $10^5\, \text{V}$ is applied between the plates, evaluate the Lorentz field. \pts{4}
\end{parts}
\hfill   \textbf{OR}

\begin{parts}
  \item Define magnetic induction $\vec{B}$, magnetization $\vec{M}$, magnetic susceptibility $\chi_m$, and permeability $\mu$. Establish the relationship between them. \pts{5}
  \item Obtain Langevin's formula for the dipolar polarizability of a polar dielectric substance. \pts{5}
  \item Evaluate the Larmor frequency of an electronic orbit when it is subjected to a magnetic induction field of $2\, \text{Tesla}$. (Given $e = 1.6  \times 10^{-19}\, \text{C}$ and $m = 9.1  \times 10^{-31}\, \text{kg}$). \pts{4}
\end{parts}

\medskip\rule{\linewidth}{0.2pt}




\noindent   \textbf{Question 4.}

\begin{parts}
  \item Distinguish between continuous and characteristic X-rays. Show the plot of X-ray intensity versus frequency observed in both cases. \pts{4}
  \item Derive Bragg's law for X-ray diffraction from a layer of atoms in a plane. \pts{6}
  \item X-rays of wavelength $0.71\, \text{\AA}$ are reflected from the (110) plane of a rock salt crystal ($d = 2.82\, \text{\AA}$). Calculate the glancing angle $ \theta$ corresponding to the second-order Bragg spectrum. \pts{4}
\end{parts}
\hfill   \textbf{OR}

\begin{parts}
  \item Define mobility of charge carriers in a semiconductor. Find the expressions for the drift current and diffusion current in a semiconductor. \pts{5}
  \item What are the various transistor biasing techniques? Draw the circuit diagram of a self-bias circuit and describe its working for stabilization. \pts{5}
  \item What is rectifier efficiency? Show that the efficiency of a half-wave rectifier is:
    \[ \eta = \frac{40.6\%}{1 + R_f / R_L} \]
    where $R_f$ is the forward resistance of the diode and $R_L$ is the load resistance. \pts{4}
\end{parts}

\medskip\rule{\linewidth}{0.2pt}




\noindent   \textbf{Question 5.}

\begin{parts}
  \item Draw the circuit diagram of a single-stage RC coupled common-emitter amplifier. Explain the function of each resistance and capacitance used in the circuit. Plot the frequency response of the amplifier. \pts{6}
  \item What is the photoelectric effect? Explain photoelectric current and retarding potential. The photoelectric threshold wavelength for a metal is $3000\, \text{\AA}$. Find the kinetic energy of an electron ejected from it by radiation of wavelength $1200\, \text{\AA}$. (Given $h = 6.62  \times 10^{-34}\, \text{J\cdot s}$ and $c = 3  \times 10^8\, \text{m/s}$). \pts{6}
  \item Explain active region, cut-off region, and saturation region in transistor operations using suitable circuit characteristics and curves. \pts{2}
\end{parts}

\vfill
\rule{\linewidth}{0.4pt}

\begin{center}\small
  \href{https://bhu-exam.vercel.app/}{Click here to visit our website}\\[4pt]
  \href{https://me-aryan.vercel.app/}{Click here to visit our contributor}\\[4pt]
  \href{https://quashan-me.vercel.app/}{Click here to visit our contributor}
\end{center}

\end{document}
